\documentclass[11pt]{article}

\usepackage[margin=2.2cm]{geometry}
\usepackage{amsmath}
\usepackage{amssymb}
\usepackage{enumitem}
\usepackage{titlesec}
\usepackage{xcolor}
\usepackage{fancyhdr}

\titleformat{\section}{\large\bfseries}{\thesection}{0.6em}{}
\titlespacing{\section}{0pt}{1.4em}{0.8em}

\pagestyle{fancy}
\fancyhf{}
\fancyhead[L]{Problems with Multiple Constraints}
\fancyhead[R]{Name: \makebox[4cm]{\hrulefill}}
\fancyfoot[C]{\thepage}
\renewcommand{\headrulewidth}{0.4pt}

\newcommand{\worksheettitle}{Solving Problems with Multiple Constraints}
\newcommand{\qspace}{\vspace{1.0cm}}

\setlist[enumerate,1]{leftmargin=1.6em, itemsep=0.4em, topsep=0.2em}
\setlist[enumerate,2]{leftmargin=1.4em, itemsep=0.3em, topsep=0.2em}

\begin{document}

\begin{center}
  {\Large\bfseries \worksheettitle}\\[0.3em]
  {\normalsize Fluency \(\cdot\) Problem Solving \(\cdot\) Reasoning \(\cdot\) Enrichment}
\end{center}

\vspace{0.3em}
\noindent\textit{Big idea: turn several written conditions into algebra, combine them into one equation, solve for the unknown, and check your answer against \textbf{every} condition.}

\vspace{0.5em}

\section{Fluency}

\subsection*{Translating words into expressions}
Write an algebraic expression for each description.

\begin{enumerate}
  \item 5 more than a number $x$.
  \qspace
  \item 4 less than twice a number $n$.
  \qspace
  \item 3 times the sum of a number $y$ and 7.
  \qspace
  \item A number $m$ decreased by 6, then doubled.
  \qspace
\end{enumerate}

\subsection*{Combining expressions}
Simplify each combination of expressions.

\begin{enumerate}[resume]
  \item If $A = x + 5$ and $B = 2x - 3$, find and simplify $A + B$.
  \qspace
  \item If $P = 3n - 2$ and $Q = n + 8$, find and simplify $2P - Q$.
  \qspace
\end{enumerate}

\subsection*{Forming and solving simple equations}
For each relationship, write an equation and solve it.

\begin{enumerate}[resume]
  \item A number increased by 8 is equal to 15. Find the number.
  \qspace
  \item Twice a number, decreased by 5, is equal to 21. Find the number.
  \qspace
\end{enumerate}

\subsection*{Solving linear equations}
Solve each equation.

\begin{enumerate}[resume]
  \item $3x + 7 = 22$
  \qspace
  \item $2(x - 4) = 10$
  \qspace
  \item $5x - 3 = 2x + 12$
  \qspace
\end{enumerate}

\subsection*{Substitution}
\begin{enumerate}[resume]
  \item If $x = 5$, find the value of $3x + 7$.
  \qspace
  \item If $y = 2x - 1$, find the value of $y$ when $x = 4$.
  \qspace
\end{enumerate}

\subsection*{Perimeter and totals}
\begin{enumerate}[resume]
  \item A rectangle has width $x$ and length $x + 3$. Write an expression for its perimeter in terms of $x$, and simplify it.
  \qspace
  \item Using your expression from the previous question, find $x$ if the perimeter is 26 cm.
  \qspace
\end{enumerate}

\newpage
\section{Problem Solving}

\noindent For each problem: define your variable, translate every constraint into algebra, combine the constraints into one equation, solve it, then answer the question that was actually asked (checking it makes sense).

\begin{enumerate}
  \item \textbf{Triangle.} In a triangle, the second side is 4 cm longer than the first side. The third side is twice the length of the first side. The perimeter of the triangle is 40 cm. Let the first side be $x$ cm.
  \begin{enumerate}
    \item Write expressions for the second and third sides in terms of $x$.
    \item Form an equation using the perimeter, and solve it to find the length of each side.
    \item Check that your three lengths could actually form a triangle.
  \end{enumerate}
  \qspace

  \item \textbf{Rectangle.} A rectangle's length is 3 cm more than twice its width. The perimeter of the rectangle is 36 cm. Let the width be $w$ cm.
  \begin{enumerate}
    \item Write an expression for the length in terms of $w$.
    \item Form an equation for the perimeter and solve it to find the width and length.
  \end{enumerate}
  \qspace

  \item \textbf{Ages.} Sarah is 3 years older than her brother Tom. In 5 years' time, the sum of their ages will be 31. Let Tom's current age be $t$ years.
  \begin{enumerate}
    \item Write an expression for Sarah's current age, and for each of their ages in 5 years' time.
    \item Form an equation for the sum of their ages in 5 years' time, and solve it to find both of their current ages.
  \end{enumerate}
  \qspace

  \item \textbf{Three numbers.} Three numbers have a sum of 100. The second number is 5 more than the first. The third number is three times the second. Let the first number be $n$. Form an equation and solve it to find all three numbers.
  \qspace

  \item \textbf{Coins.} A charity box contains only \$1 and \$2 coins. There are 5 more \$2 coins than \$1 coins. The total value of all the coins is \$70. Let the number of \$1 coins be $x$.
  \begin{enumerate}
    \item Write an expression for the number of \$2 coins, and for the total value of all the coins.
    \item Form an equation and solve it to find how many of each coin there are.
  \end{enumerate}
  \qspace

  \item \textbf{Ribbon.} A 200 cm ribbon is cut into three pieces. The second piece is twice as long as the first piece. The third piece is 20 cm longer than the second piece. Let the first piece be $x$ cm. Form an equation and solve it to find the length of each piece.
  \qspace
\end{enumerate}

\newpage
\section{Reasoning}

\begin{enumerate}
  \item \textbf{Find the error.} A student was asked: \textit{``The width of a rectangle is $w$. The length is 6 cm more than the width. The perimeter is 28 cm. Find the width.''} Their working was:
  \[
    w + (w + 6) = 28, \qquad 2w + 6 = 28, \qquad 2w = 22, \qquad w = 11.
  \]
  Identify the error in the student's working, explain why it is incorrect, and give the correct solution.
  \qspace

  \item \textbf{Necessity of information.} Consider the problem: \textit{``The length of a rectangle is 4 cm more than its width, and the perimeter is 32 cm. Find the dimensions.''} Explain why \emph{both} pieces of information --- the relationship between length and width, \emph{and} the perimeter --- are necessary to find exact dimensions. What could you find (and what could you not find) if you were only given one of these facts?
  \qspace

  \item \textbf{Comparing methods.} Two numbers differ by 7. Their sum is 41.
  \begin{itemize}
    \item Student A let $x$ be the smaller number and wrote the equation $x + (x + 7) = 41$.
    \item Student B let $x$ be the larger number and $y$ be the smaller number, and wrote the two equations $x - y = 7$ and $x + y = 41$.
  \end{itemize}
  Solve the problem using \emph{both} methods. Then explain one advantage and one disadvantage of each approach.
  \qspace

  \item \textbf{Always, sometimes, or never true.} Consider the statement: \textit{``If a rectangle's length is twice its width, then its perimeter will always be a multiple of 6.''} Decide whether this statement is always, sometimes, or never true. Justify your answer using algebra and/or examples.
  \qspace

  \item \textbf{Justifying a step.} In solving the equation $3(x + 4) = 2x + 19$, a student's first step was to write $3x + 12 = 2x + 19$.
  \begin{enumerate}
    \item Justify this step using a named algebraic law.
    \item Complete the solution, and verify that your value of $x$ satisfies the \emph{original} equation.
  \end{enumerate}
  \qspace

  \item \textbf{Changing a constraint.} In the Ribbon problem from Problem Solving (Question 6), the third piece was described as 20 cm \emph{longer} than the second piece. Suppose instead it was 20 cm \emph{shorter} than the second piece, with the total length of ribbon unchanged. Re-solve the problem under this new constraint. Does this change still give a sensible answer? Explain what feature of an equation's solution would warn you if it did not.
  \qspace

  \item \textbf{Checking validity and restrictions.} A problem states: \textit{``A rectangle has width $w$ and length $2w - 10$. The perimeter is 10 cm.''}
  \begin{enumerate}
    \item Solve for $w$ algebraically.
    \item Explain why your algebraically correct solution is \emph{not} a valid rectangle.
    \item What restriction should be placed on $w$ for this problem to describe a genuine rectangle, and why?
  \end{enumerate}
  \qspace
\end{enumerate}

\newpage
\section{Enrichment}

\begin{enumerate}
  \item \textbf{Generalising with a parameter.} A rectangle has width $w$ and length equal to $k$ times its width, where $k > 1$ is a constant. Its perimeter is $P$.
  \begin{enumerate}
    \item Find an expression for $w$ in terms of $k$ and $P$.
    \item Hence find an expression for the length in terms of $k$ and $P$.
    \item Use your formulas to find the width and length when $k = 3$ and $P = 48$. Check that these values satisfy the original conditions.
    \item Describe what happens to the \emph{shape} of the rectangle as $k$ increases, for a fixed perimeter $P$. Does it become more square or more elongated? Justify your answer using your formula for $w$.
  \end{enumerate}
  \qspace

  \item \textbf{Reverse-engineering a problem.} The answer to a certain ``three consecutive even numbers'' problem is that the numbers are 24, 26 and 28.
  \begin{enumerate}
    \item Write your own word problem, involving at least two constraints and in a similar style to the Problem Solving section, for which $x = 24$ (the smallest number) is the correct solution.
    \item Solve your own problem algebraically and show that it does give this answer.
  \end{enumerate}
  \qspace

  \item \textbf{Investigating all possible values.} A triangle has sides of length $x$, $x + 4$, and $2x - 1$ (all in centimetres).
  \begin{enumerate}
    \item Using the triangle inequality (the sum of any two sides must exceed the third side), write down all the inequalities that $x$ must satisfy.
    \item Hence find the full range of values of $x$ for which a valid triangle exists.
    \item If $x$ must be a whole number, list every possible value of $x$, together with the perimeter of the corresponding triangle.
  \end{enumerate}
  \qspace

  \item \textbf{Prove or disprove.} A student claims: \textit{``For any triangle with sides $x$, $x+a$, and $x+b$, where $a$ and $b$ are positive constants with $a \neq b$, it is always possible to choose a value of $x$ that makes a valid triangle.''}

  Prove or disprove this claim using algebra. If you disprove it, give a counterexample. If you prove it, show precisely why a suitable value of $x$ must always exist.
  \qspace
\end{enumerate}

\end{document}
